\chapter{Átomos polieletrônicos e símbolos de termo}\label{ch:b3:many-electron-atoms}

Em 1868, uma raia amarela no espectro do Sol revelou um elemento
desconhecido na Terra, o hélio. Quando ele foi finalmente isolado em
laboratório, em 1895, seu espectro parecia o de \emph{dois} elementos: duas
famílias de raias, atribuídas por algum tempo ao ``orto-hélio'' e ao
``para-hélio'', que nunca trocavam luz entre si. Há um só hélio. As duas
famílias são os \omterm{def:b3:many-electron-atoms:singlet-triplet}{estados tripleto} e singleto do mesmo átomo, mantidos
separados pelo spin do elétron e por uma regra mais fundamental que
qualquer força: a função de onda de dois elétrons deve mudar de sinal quando
eles são permutados. Este capítulo segue essa regra desde o princípio de
Pauli até os \omterm{def:b3:many-electron-atoms:term}{símbolos de termo} com que os espectroscopistas rotulam cada
estado de cada átomo e íon --- inclusive os íons de metais de transição
cujas cores são explicadas no \cref{ch:b3:complex-spectra-magnetism}.

\begin{recall}
O volume do 1.º ano rotulou os elétrons por quatro números quânticos $n$, $l$,
$m_l$, $m_s$, construiu configurações fundamentais com o princípio de Pauli,
a regra de Klechkowski e a regra de Hund, e contou elétrons desemparelhados.
O volume do 2.º ano deu as energias orbitais e as regras de Slater para a
blindagem. O \cref{ch:b3:quantum-model-systems} forneceu os \omterm{def:b3:quantum-model-systems:operator}{operadores}, os
\omterm{def:b3:quantum-model-systems:operator}{autovalores}, o momento angular do rotor ($\hat L^2$ com \omterm{def:b3:quantum-model-systems:operator}{autovalores}
$l(l+1)\hbar^2$) e o átomo de hidrogênio.
\end{recall}

\section{Spin e spin-orbitais}

O elétron possui um momento angular intrínseco, o spin, de número quântico
$s = \frac12$: $\hat S^2$ tem o único \omterm{def:b3:quantum-model-systems:operator}{autovalor} $s(s+1)\hbar^2 =
\frac34\hbar^2$, e $\hat S_z$ os dois \omterm{def:b3:quantum-model-systems:operator}{autovalores} $m_s\hbar = \pm\frac12\hbar$.
As duas funções de spin são escritas $\alpha$ ($m_s = +\frac12$) e $\beta$
($m_s = -\frac12$); elas são ortonormais, $\langle\alpha|\alpha\rangle =
\langle\beta|\beta\rangle = 1$ e $\langle\alpha|\beta\rangle = 0$. O spin não tem
equivalente clássico; não é a rotação de uma pequena esfera e nem sequer
entra em um hamiltoniano não relativístico.

\begin{definition}[Spin-orbital]\label{def:b3:many-electron-atoms:spin-orbital}
Um \emph{spin-orbital}\index{spin-orbital} é o produto de um orbital espacial
$\phi(\mathbf r)$ por uma função de spin: $\phi\alpha$ ou $\phi\beta$. Um
spin-orbital contém a descrição completa de um elétron.
\end{definition}

Um orbital atômico rotulado $(n, l, m_l)$ no volume do 1.º ano dá, assim, dois
\omterm{def:b3:many-electron-atoms:spin-orbital}{spin-orbitais}, e é por isso que uma subcamada comporta $2(2l+1)$ elétrons. A
verdadeira questão é por que um \omterm{def:b3:many-electron-atoms:spin-orbital}{spin-orbital} comporta no máximo um.

\section{Elétrons indistinguíveis e o determinante de Slater}

\begin{definition}[Partículas indistinguíveis, função de onda antissimétrica]\label{def:b3:many-electron-atoms:indistinguishable}
Partículas são \emph{indistinguíveis}\index{particulas indistinguiveis@partículas indistinguíveis} se
nenhuma medida consegue diferenciá-las: todo elétron é idêntico a qualquer
outro. Uma função de onda de vários elétrons é uma \emph{função de onda
antissimétrica}\index{funcao de onda antissimetrica@função de onda antissimétrica} se muda de sinal quando
as coordenadas (espaciais e de spin) de dois elétrons quaisquer são permutadas:
$\Psi(\dots,i,\dots,j,\dots) = -\Psi(\dots,j,\dots,i,\dots)$.
\end{definition}

A indistinguibilidade, por si só, exige que $|\Psi|^2$ não mude com uma
permutação, de modo que $\Psi$ pode, no máximo, mudar de sinal. A natureza
escolheu: a função de onda de qualquer sistema de elétrons --- de quaisquer
férmions --- é antissimétrica. É um postulado da mecânica quântica,
confirmado por todo espectro atômico.

\begin{definition}[Determinante de Slater]\label{def:b3:many-electron-atoms:slater}
Para $N$ elétrons em $N$ \omterm{def:b3:many-electron-atoms:spin-orbital}{spin-orbitais} diferentes $\chi_1, \dots, \chi_N$, o
\emph{determinante de Slater}\index{determinante de Slater} é
\[
\Psi = \frac{1}{\sqrt{N!}}\begin{vmatrix}
\chi_1(1) & \chi_2(1) & \cdots & \chi_N(1)\\
\chi_1(2) & \chi_2(2) & \cdots & \chi_N(2)\\
\vdots & & & \vdots\\
\chi_1(N) & \chi_2(N) & \cdots & \chi_N(N)
\end{vmatrix},
\]
onde a linha $i$ contém o elétron $i$ em cada \omterm{def:b3:many-electron-atoms:spin-orbital}{spin-orbital}. Ele é escrito, de
forma abreviada, $|\chi_1\chi_2\dots\chi_N|$.
\end{definition}

\begin{theorem}[O princípio de Pauli]\label{thm:b3:many-electron-atoms:pauli}
Um \omterm{def:b3:many-electron-atoms:slater}{determinante de Slater} é antissimétrico e se anula se dois elétrons
ocupam o mesmo \omterm{def:b3:many-electron-atoms:spin-orbital}{spin-orbital}. Portanto, dois elétrons de um átomo não podem
ter os mesmos quatro números quânticos.
\end{theorem}

\begin{proof}
Permutar os elétrons $i$ e $j$ permuta duas linhas do determinante, o que
muda seu sinal: a antissimetria vale para todo par. Se dois \omterm{def:b3:many-electron-atoms:spin-orbital}{spin-orbitais}
são iguais, duas colunas são iguais e o determinante é nulo: tal estado não
existe. Dois elétrons do mesmo átomo com os mesmos $n$, $l$, $m_l$ e $m_s$
estariam no mesmo \omterm{def:b3:many-electron-atoms:spin-orbital}{spin-orbital}.
\end{proof}

\begin{example}[O estado fundamental do hélio]\label{ex:b3:many-electron-atoms:helium-ground}
Com os dois elétrons em $1s$,
\[
\Psi = \frac{1}{\sqrt2}\begin{vmatrix}1s\alpha(1) & 1s\beta(1)\\ 1s\alpha(2) &
1s\beta(2)\end{vmatrix} = 1s(1)\,1s(2)\cdot\frac{1}{\sqrt2}\big[\alpha(1)\beta(2)
- \beta(1)\alpha(2)\big].
\]
A parte espacial é simétrica, a parte de spin antissimétrica: os dois spins
estão emparelhados, com spin total nulo.
\end{example}

\section{Hélio: repulsão coulombiana e troca}

Excite um elétron do hélio para $2s$. Quatro determinantes podem ser
construídos a partir de $1s\alpha$, $1s\beta$, $2s\alpha$, $2s\beta$ com um
elétron em cada orbital. Rearranjados, eles se tornam produtos de uma função
espacial por uma função de spin:
\[
\Psi_\pm = \frac{1}{\sqrt2}\big[1s(1)2s(2) \pm 2s(1)1s(2)\big]\times(\text{função
de spin}).
\]
A função espacial antissimétrica $\Psi_-$ deve ser associada a uma função de
spin simétrica, das quais há três: $\alpha(1)\alpha(2)$,
$\beta(1)\beta(2)$ e $\frac{1}{\sqrt2}[\alpha(1)\beta(2) + \beta(1)\alpha(2)]$, as
três componentes $M_S = 1, 0, -1$ de um spin total $S = 1$. A função simétrica
$\Psi_+$ deve ser associada à única função de spin antissimétrica, $S = 0$.

\begin{definition}[Estados singleto e tripleto]\label{def:b3:many-electron-atoms:singlet-triplet}
Um estado de spin total $S = 0$ é um \emph{estado singleto}\index{estado singleto};
um estado de spin total $S = 1$, cujas três componentes $M_S = 1, 0, -1$ têm a
mesma energia na ausência de campo magnético, é um \emph{estado
tripleto}\index{estado tripleto}.
\end{definition}

\begin{definition}[Integral de troca]\label{def:b3:many-electron-atoms:exchange}
Para dois orbitais $a$ e $b$ e a repulsão eletrônica $e^2/4\pi\varepsilon_0
r_{12}$, a \emph{integral de troca}\index{integral de troca} é
\[
K_{ab} = \iint a(1)b(2)\,\frac{e^2}{4\pi\varepsilon_0r_{12}}\,b(1)a(2)\,\dd\tau_1\dd\tau_2,
\]
a mesma integral que a repulsão eletrônica clássica $J_{ab}$ (com
$a(1)b(2)$ dos dois lados), exceto que os elétrons estão trocados de um
lado. $K_{ab}$ é positiva e não tem equivalente clássico.
\end{definition}

\begin{proposition}[Energias do singleto e do tripleto]\label{prop:b3:many-electron-atoms:helium}
Em primeira ordem na repulsão eletrônica, os estados $1s2s$ do hélio têm
energias $E_\pm = E_{1s} + E_{2s} + J \pm K$, o singleto ($+$) acima do
tripleto ($-$) por $2K$.
\end{proposition}

\begin{proof}
A energia não perturbada de ambos os $\Psi_\pm$ é $E_{1s} + E_{2s}$ (orbitais
hidrogenoides de carga nuclear 2). A correção de primeira ordem é o \omterm{def:b3:quantum-model-systems:hermitian}{valor
esperado} da repulsão $\hat V_{12}$ no estado normalizado:
$\frac12\langle 1s(1)2s(2) \pm 2s(1)1s(2)|\hat V_{12}|1s(1)2s(2) \pm
2s(1)1s(2)\rangle$. Os dois termos diretos dão $\frac12(J + J)$; os dois termos
cruzados dão $\pm\frac12(K + K)$, pois $\hat V_{12}$ é simétrico nos
elétrons. Logo $\Delta E = J \pm K$, e as funções de spin, normalizadas e
não afetadas por $\hat V_{12}$, apenas multiplicam por 1.
\end{proof}

O tripleto fica mais abaixo, embora o hamiltoniano não contenha o spin: sua
função espacial se anula quando $\mathbf r_1 = \mathbf r_2$, de modo que os
dois elétrons se evitam e se repelem menos. Esse ``buraco de Fermi'' é o
conteúdo físico da primeira regra de Hund.

\begin{example}[A integral de troca do hélio, medida]\label{ex:b3:many-electron-atoms:K}
% ledger: lev:He.2s3S1, lev:He.2s1S0
Os níveis $1s2s$ do hélio ficam a \qty{159856}{cm^{-1}} (${}^3$S) e
\qty{166277}{cm^{-1}} (${}^1$S) acima do estado fundamental: o singleto está
\qty{6421}{cm^{-1}}, \qty{0.796}{eV}, mais alto, e $K(1s,2s) = \qty{0.398}{eV}$.
A figura abaixo mostra as duas famílias.
\end{example}

\section{Acoplamento de Russell--Saunders e símbolos de termo}

Em um átomo leve, os momentos angulares orbitais dos elétrons se somam em
um total $\mathbf L$, seus spins em um total $\mathbf S$, e só então
$\mathbf L$ e $\mathbf S$ se acoplam fracamente entre si.

\begin{definition}[Acoplamento de Russell--Saunders]\label{def:b3:many-electron-atoms:russell-saunders}
No \emph{acoplamento de Russell--Saunders}\index{acoplamento de Russell--Saunders}, os
estados de uma configuração são rotulados pelo \emph{número quântico de momento
angular orbital total}\index{numero quantico de momento angular orbital total@número quântico de momento angular orbital total}
$L$ ($\hat{\mathbf L}^2 = L(L+1)\hbar^2$, com $M_L = \sum m_l$), pelo \emph{número
quântico de spin total}\index{numero quantico de spin total@número quântico de spin total} $S$ ($M_S = \sum m_s$)
e pelo \emph{número quântico de momento angular total}\index{numero quantico de momento angular total@número quântico de momento angular total}
$J$, que assume os valores $|L - S|, \dots, L + S$.
\end{definition}

\begin{definition}[Termo espectroscópico, multiplicidade de spin, símbolo de termo]\label{def:b3:many-electron-atoms:term}
Um \emph{termo espectroscópico}\index{termo espectroscopico@termo espectroscópico} é o conjunto dos
$(2L+1)(2S+1)$ estados de uma configuração com $L$ e $S$ dados. Sua
\emph{multiplicidade de spin}\index{multiplicidade de spin} é $2S + 1$. O \emph{símbolo de termo}\index{simbolo de termo@símbolo de termo}
é ${}^{2S+1}L$, com as letras S, P, D, F, G, H
para $L = 0, 1, 2, 3, 4, 5$; um estado de $J$ dado é escrito ${}^{2S+1}L_J$ ---
por exemplo, \termsym{3}{P}{2}.
\end{definition}

\begin{proposition}[Camadas completas]\label{prop:b3:many-electron-atoms:closed-shell}
Uma subcamada completa tem $L = 0$ e $S = 0$; ela não contribui para o
\omterm{def:b3:many-electron-atoms:term}{símbolo de termo}.
\end{proposition}

\begin{proof}
Em uma subcamada completa, cada $m_l$ aparece duas vezes e cada $m_s = \pm\frac12$
aparece $2l+1$ vezes, logo $M_L = 0$ e $M_S = 0$; e só há uma maneira de
preenchê-la (um único determinante). Um único estado com $M_L = M_S = 0$ só pode
pertencer a $L = 0$, $S = 0$.
\end{proof}

\begin{proposition}[Contagem dos estados de uma configuração]\label{prop:b3:many-electron-atoms:count}
$N$ elétrons em uma subcamada de $2(2l+1)$ \omterm{def:b3:many-electron-atoms:spin-orbital}{spin-orbitais} dão
$\binom{2(2l+1)}{N}$ determinantes, e as \omterm{def:b3:quantum-model-systems:degenerate}{degenerescências} $(2L+1)(2S+1)$ de
seus termos somam esse número.
\end{proposition}

\begin{proof}
Um determinante é fixado pelo conjunto dos \omterm{def:b3:many-electron-atoms:spin-orbital}{spin-orbitais} ocupados (sua ordem
só muda o sinal): escolhem-se $N$ entre os $2(2l+1)$. Os termos são os
mesmos estados reagrupados, logo as contagens devem coincidir.
\end{proof}

\begin{method}[Termos de uma configuração]\label{met:b3:many-electron-atoms:terms}
\begin{enumerate}
  \item Liste os determinantes permitidos pelo princípio de Pauli e
        organize-os em uma tabela por $M_L$ e $M_S$.
  \item Encontre o maior $M_L$; com o maior $M_S$ que o acompanha, ele
        inicia um termo com $L = M_L$, $S = M_S$.
  \item Risque um determinante para cada par $(M_L, M_S)$ com
        $|M_L| \le L$, $|M_S| \le S$.
  \item Repita com o que sobra, até esvaziar a tabela; confira a contagem.
\end{enumerate}
\end{method}

\begin{example}[Os termos de $p^2$]\label{ex:b3:many-electron-atoms:p2}
Dois elétrons em $p$ ($m_l = 1, 0, -1$) dão $\binom62 = 15$ determinantes.
Sua tabela por $M_L$ e $M_S$ é:
\begin{center}\small
\begin{tabular}{c|ccc}
\hline
 & $M_S = 1$ & $M_S = 0$ & $M_S = -1$\\ \hline
$M_L = 2$ & -- & $(1^+,1^-)$ & --\\
$M_L = 1$ & $(1^+,0^+)$ & $(1^+,0^-)$, $(1^-,0^+)$ & $(1^-,0^-)$\\
$M_L = 0$ & $(1^+,-1^+)$ & $(1^+,-1^-)$, $(1^-,-1^+)$, $(0^+,0^-)$ & $(1^-,-1^-)$\\
$M_L = -1$ & $(0^+,-1^+)$ & $(0^+,-1^-)$, $(0^-,-1^+)$ & $(0^-,-1^-)$\\
$M_L = -2$ & -- & $(-1^+,-1^-)$ & --\\ \hline
\end{tabular}
\end{center}
(com $m_l^\pm$ para $m_s = \pm\frac12$). $M_L = 2$ só ocorre com $M_S = 0$:
um termo ${}^1$D (5 estados). O maior $M_L = 1$ restante vem com
$M_S = 1$: um termo ${}^3$P (9 estados). Resta um estado com $M_L = M_S = 0$: um
termo ${}^1$S. De fato, $5 + 9 + 1 = 15$.
\end{example}

\begin{example}[O carbono, medido]\label{ex:b3:many-electron-atoms:carbon}
% ledger: lev:C.3P1, lev:C.3P2, lev:C.1D2, lev:C.1S0
A configuração fundamental $2p^2$ do carbono dá os níveis \termsym{3}{P}{0},
\termsym{3}{P}{1}, \termsym{3}{P}{2} em 0, 16.4 e \qty{43.4}{cm^{-1}}, depois
\termsym{1}{D}{2} em \qty{10193}{cm^{-1}} e \termsym{1}{S}{0} em
\qty{21648}{cm^{-1}}: o termo ${}^3$P é o mais baixo, como preveem as regras de Hund.
\end{example}

\begin{omfigure}
\begin{tikzpicture}
\begin{axis}[name=main, width=0.62\linewidth, height=6.4cm, xmin=-0.3, xmax=6.6,
  ymin=-1500, ymax=23500, axis y line=left, axis x line=none,
  ylabel={energia / cm$^{-1}$}, unbounded coords=jump, clip=false,
  scaled y ticks=false, ytick={0,10000,20000},
  tick label style={font=\small}, label style={font=\small}]
\addplot[omstate] table[x=x, y=E] {figdata/out/bachelor-3-atomic-levels-carbon.dat};
\draw[densely dotted, gray] (axis cs:1,4858.5) -- (axis cs:2,29.6);
\draw[densely dotted, gray] (axis cs:1,4858.5) -- (axis cs:2,10192.7);
\draw[densely dotted, gray] (axis cs:1,4858.5) -- (axis cs:2,21648);
\draw[densely dotted, gray] (axis cs:3,10192.7) -- (axis cs:4,10192.7);
\draw[densely dotted, gray] (axis cs:3,21648) -- (axis cs:4,21648);
\node[font=\small, anchor=south] at (axis cs:0.5,4858.5) {$2p^2$};
\node[font=\small, anchor=south] at (axis cs:2.5,21648) {${}^1$S};
\node[font=\small, anchor=south] at (axis cs:2.5,10192.7) {${}^1$D};
\node[font=\small, anchor=south] at (axis cs:2.5,29.6) {${}^3$P};
\node[font=\small, anchor=west] at (axis cs:5.05,21648) {\termsym{1}{S}{0}};
\node[font=\small, anchor=west] at (axis cs:5.05,10192.7) {\termsym{1}{D}{2}};
\node[font=\scriptsize, anchor=north] at (axis cs:0.5,-600) {configuração};
\node[font=\scriptsize, anchor=north] at (axis cs:2.5,-600) {termos};
\node[font=\scriptsize, anchor=north] at (axis cs:4.5,-600) {níveis};
\draw[gray, ->] (axis cs:3.05,60) -- (axis cs:6.9,2500);
\end{axis}
\begin{axis}[at={(main.south east)}, anchor=south west, xshift=1.2cm, yshift=0.9cm,
  width=3.6cm, height=3.6cm, xmin=-0.1, xmax=1.1, ymin=-5, ymax=50,
  axis y line=right, axis x line=none, unbounded coords=jump, ytick={0,16.4,43.4},
  yticklabels={0,16.4,43.4}, tick label style={font=\scriptsize}, clip=false,
  title={\scriptsize ${}^3$P, ampliado}, title style={yshift=-4pt}]
\addplot[omstate] table[x=x, y=E] {figdata/out/bachelor-3-atomic-levels-carbon-zoom.dat};
\node[font=\scriptsize, anchor=east] at (axis cs:0,0) {$J = 0$};
\node[font=\scriptsize, anchor=east] at (axis cs:0,16.4) {$J = 1$};
\node[font=\scriptsize, anchor=east] at (axis cs:0,43.4) {$J = 2$};
\end{axis}
\end{tikzpicture}

{\small A configuração fundamental do carbono: uma configuração, três termos
(em suas energias medidas; o termo ${}^3$P na média ponderada de seus
níveis, a configuração na média dos quinze estados) e cinco
níveis. O desdobramento spin--órbita de ${}^3$P (ampliado, em
\unit{cm^{-1}}) é algumas centenas de vezes menor que as separações entre
termos.}
\end{omfigure}

\section{Regras de Hund e acoplamento spin--órbita}

\begin{proposition}[Regras de Hund para os termos]\label{prop:b3:many-electron-atoms:hund-terms}
Para a configuração fundamental de um átomo ou íon, o termo mais baixo é o
que tem (1) o maior $S$ e depois (2), para esse $S$, o maior $L$. Seu nível
mais baixo tem (3) $J = |L - S|$ se a subcamada está menos que semipreenchida, $J = L +
S$ se está mais que semipreenchida.
\end{proposition}
\begin{proof}[Estatuto]
Essas regras são estabelecidas experimentalmente, apenas para as
configurações fundamentais; não são teoremas. A regra 1 é a estabilização
por troca da seção anterior; a regra 2 diz que elétrons que giram no mesmo
sentido se encontram menos; a regra 3 decorre do sinal do \omterm{def:b3:many-electron-atoms:spin-orbit}{acoplamento
spin--órbita}, mais adiante.
\end{proof}

\begin{method}[O termo fundamental a partir de um diagrama de caixas]\label{met:b3:many-electron-atoms:ground-term}
\begin{enumerate}
  \item Preencha as caixas da subcamada aberta a partir de $m_l = +l$ para baixo,
        um elétron por caixa com spins paralelos, depois emparelhe a partir de $+l$ de novo.
  \item $S = \frac12 \times$(número de elétrons desemparelhados); $L = |\sum m_l|$.
  \item $J = |L - S|$ se a subcamada está menos que semipreenchida, $L + S$ se
        está mais que semipreenchida; exatamente semipreenchida, $L = 0$ e $J = S$.
\end{enumerate}
\end{method}

\begin{example}[Termos fundamentais]\label{ex:b3:many-electron-atoms:ground-terms}
% ledger: lev:Ti2.3P0, lev:Ni2.3P0, lev:Cr3.4P1_2
\ce{Ti^2+}, $d^2$: caixas $m_l = 2, 1$ preenchidas, $S = 1$, $L = 3$, $J = 2$:
\termsym{3}{F}{2}. \ce{Cr^3+}, $d^3$: $S = \frac32$, $L = 2 + 1 + 0 = 3$,
\termsym{4}{F}{3/2}. \ce{Fe^2+}, $d^6$: cinco para cima, um para baixo em $m_l = 2$:
$S = 2$, $L = 2$, mais que semipreenchida, \termsym{5}{D}{4}. \ce{Ni^2+}, $d^8$:
$S = 1$, $L = 3$, \termsym{3}{F}{4}. Cada um concorda com o nível fundamental
que a espectroscopia atômica mede.
\end{example}

\begin{definition}[Acoplamento spin--órbita, estrutura fina]\label{def:b3:many-electron-atoms:spin-orbit}
O \emph{acoplamento spin--órbita}\index{acoplamento spin--orbita@acoplamento spin--órbita} é a interação
entre os momentos magnéticos do spin de um elétron e de seu movimento
orbital, escrita $A\,\hat{\mathbf L}\cdot\hat{\mathbf S}/\hbar^2$ para um termo. Ele
desdobra um termo em seus níveis de $J$ diferentes: a \emph{estrutura
fina}\index{estrutura fina} do termo.
\end{definition}

\begin{proposition}[Regra dos intervalos de Landé]\label{prop:b3:many-electron-atoms:lande}
Dentro de um termo, $E(J) = E_0 + \frac{A}{2}[J(J+1) - L(L+1) - S(S+1)]$, de modo que
$E(J) - E(J - 1) = AJ$. $A > 0$ (normal) para uma subcamada menos que
semipreenchida, $A < 0$ (invertida) para uma mais que semipreenchida.
\end{proposition}

\begin{proof}
$\hat{\mathbf J} = \hat{\mathbf L} + \hat{\mathbf S}$ dá $\hat{\mathbf J}^2 =
\hat{\mathbf L}^2 + \hat{\mathbf S}^2 + 2\hat{\mathbf L}\cdot\hat{\mathbf S}$, logo
$\hat{\mathbf L}\cdot\hat{\mathbf S} = \frac12(\hat{\mathbf J}^2 - \hat{\mathbf L}^2 -
\hat{\mathbf S}^2)$, cujo \omterm{def:b3:quantum-model-systems:operator}{autovalor} em um estado de $J$, $L$, $S$ dados é
$\frac{\hbar^2}{2}[J(J+1) - L(L+1) - S(S+1)]$. A diferença entre $J$ e
$J - 1$ é $\frac A2[J(J+1) - (J-1)J] = AJ$. O sinal de $A$ é admitido: uma
subcamada mais que semipreenchida se comporta como o número correspondente
de buracos positivos.
\end{proof}

\begin{example}[Testando a regra dos intervalos]\label{ex:b3:many-electron-atoms:lande-test}
% ledger: lev:C.3P1, lev:C.3P2, lev:O.3P1, lev:O.3P0
Para ${}^3$P, a regra prevê intervalos na razão $J = 2 : 1$, isto é, 2.
Carbono: $(43.4 - 16.4)/16.4 = 1.65$. Oxigênio ($2p^4$, invertido, \termsym{3}{P}{2}
mais baixo, depois $J = 1$ a \qty{158.3}{cm^{-1}} e $J = 0$ a \qty{227.0}{cm^{-1}}):
$158.3/68.7 = 2.30$. A regra vale com 20\,\% de margem: o \omterm{def:b3:many-electron-atoms:russell-saunders}{acoplamento de
Russell--Saunders} é uma boa descrição dos átomos leves, não uma descrição exata.
\end{example}

\begin{omfigure}
\begin{tikzpicture}[font=\small]
  \draw[omstate] (0,0) -- (1.6,0) node[midway, below] {$3s\ {}^2$S$_{1/2}$};
  \draw[omstate] (0,3.0) -- (1.6,3.0);
  \node[anchor=east] at (-0.1,3.0) {$3p$};
  \draw[omstate] (2.8,2.75) -- (4.4,2.75) node[right] {${}^2$P$_{1/2}$};
  \draw[omstate] (2.8,3.25) -- (4.4,3.25) node[right] {${}^2$P$_{3/2}$};
  \draw[densely dotted, gray] (1.6,3.0) -- (2.8,2.75) (1.6,3.0) -- (2.8,3.25);
  \draw[omfluo] (3.3,2.73) -- (3.3,0.02) node[pos=0.55, left] {D$_1$};
  \draw[omfluo] (3.9,3.23) -- (3.9,0.02) node[pos=0.55, right] {D$_2$};
  \draw[omstate] (2.8,0) -- (4.4,0);
  \draw[densely dotted, gray] (1.6,0) -- (2.8,0);
  \draw[decorate, decoration={brace, amplitude=3pt}] (5.5,3.27) -- (5.5,2.73)
     node[midway, right=3pt, align=left] {\footnotesize \qty{17.2}{cm^{-1}}\\[-1pt]\footnotesize (fora de escala)};
\end{tikzpicture}

{\small As raias D do sódio. O nível $3p$ é desdobrado pelo \omterm{def:b3:many-electron-atoms:spin-orbit}{acoplamento
spin--órbita} em ${}^2$P$_{1/2}$ e ${}^2$P$_{3/2}$; a emissão para o nível
fundamental $3s$ dá duas raias amarelas, D$_1$ em \qty{589.76}{nm} e D$_2$ em
\qty{589.16}{nm} (comprimentos de onda no vácuo).}
\end{omfigure}

\begin{example}[O dubleto do sódio]\label{ex:b3:many-electron-atoms:sodium}
% ledger: lev:Na.3p1_2, lev:Na.3p3_2
Os níveis $3p$ do sódio ficam em 16956.17 e \qty{16973.37}{cm^{-1}}: as raias D
estão em $10^7/16956.17 = \qty{589.76}{nm}$ e \qty{589.16}{nm} no vácuo,
separadas por \qty{17.20}{cm^{-1}}. Para ${}^2$P, $E(\frac32) - E(\frac12) =
\frac32A$, logo $A = \qty{11.47}{cm^{-1}}$.
\end{example}

\begin{proposition}[Regras de seleção para os átomos]\label{prop:b3:many-electron-atoms:selection-atoms}
No \omterm{def:b3:many-electron-atoms:russell-saunders}{acoplamento de Russell--Saunders}, uma transição de dipolo elétrico exige
$\Delta S = 0$, $\Delta L = 0, \pm1$, $\Delta J = 0, \pm1$ (mas não $J = 0 \to J
= 0$) e uma mudança de paridade: para o salto de um único elétron, $\Delta l = \pm1$.
\end{proposition}
\begin{proof}[Estatuto]
$\Delta S = 0$ é demonstrada no \cref{ch:b3:electronic-spectroscopy}: o \omterm{def:b3:quantum-model-systems:operator}{operador}
dipolo não age sobre o spin. As demais decorrem do caráter vetorial do
\omterm{def:b3:quantum-model-systems:operator}{operador} dipolo e são admitidas aqui.
\end{proof}

\begin{omfigure}
\begin{tikzpicture}
\begin{axis}[width=0.62\linewidth, height=6.2cm, xmin=-0.5, xmax=3.5,
  ymin=158, ymax=173, axis y line=left, axis x line=none,
  ylabel={energia / $10^3$ cm$^{-1}$}, unbounded coords=jump, clip=false,
  tick label style={font=\small}, label style={font=\small}]
\addplot[omstate] table[x=x, y=E] {figdata/out/bachelor-3-atomic-levels-helium.dat};
\node[font=\small, anchor=west] at (axis cs:1.05,166.277) {$1s2s\ {}^1$S};
\node[font=\small, anchor=west] at (axis cs:1.05,171.135) {$1s2p\ {}^1$P};
\node[font=\small, anchor=west] at (axis cs:3.05,159.856) {$1s2s\ {}^3$S};
\node[font=\small, anchor=west] at (axis cs:3.05,169.087) {$1s2p\ {}^3$P};
\node[font=\small, anchor=south] at (axis cs:0.5,172.3) {singletos};
\node[font=\small, anchor=south] at (axis cs:2.5,172.3) {tripletos};
\draw[omfluo] (axis cs:0.5,171.0) -- (axis cs:0.5,166.42) node[pos=0.5, left, font=\scriptsize] {\qty{2059}{nm}};
\draw[omfluo] (axis cs:2.5,168.95) -- (axis cs:2.5,160.0) node[pos=0.5, left, font=\scriptsize] {\qty{1083}{nm}};
\draw[densely dotted, gray] (axis cs:-0.3,159.856) -- (axis cs:2,159.856);
\draw[<->, omThm] (axis cs:-0.2,159.856) -- (axis cs:-0.2,166.277) node[pos=0.5, right, font=\scriptsize] {$2K$};
\end{axis}
\end{tikzpicture}

{\small Os dois ``hélios''. Níveis singleto e tripleto das configurações
$1s2s$ e $1s2p$ em suas energias medidas acima do estado fundamental (o
termo ${}^3$P na média de seus níveis). As raias ligam apenas estados da
mesma família; cada tripleto fica abaixo de seu singleto, por $2K$.}
\end{omfigure}

\begin{history}[Dois hélios e o princípio de exclusão]
O hélio recebeu o nome do Sol em 1868 e foi encontrado na Terra por William
Ramsay em 1895. Suas duas séries de raias intrigaram os espectroscopistas
por trinta anos. Em 1925, Wolfgang Pauli propôs que dois elétrons nunca
compartilham os mesmos números quânticos, no mesmo ano em que o spin do
elétron foi proposto por George Uhlenbeck e Samuel Goudsmit; em 1926, Werner
Heisenberg mostrou que a simetria da função de onda basta para separar o
hélio em suas famílias singleto e tripleto.
\end{history}

\begin{inthelab}[Resolver o dubleto do sódio]
Uma lâmpada de sódio é colocada diante da fenda de um espectrômetro de rede.
Com uma rede de 600 linhas por milímetro, os ângulos de primeira ordem das
duas raias diferem de cerca de \num{0.02}\textdegree: um bom espectrômetro as
separa, um prisma de aula não. A lâmpada funciona quente; só é manipulada
depois de esfriar.
\end{inthelab}

\section{Exercícios}

\begin{exercise}[$\star$]\label{exo:b3:many-electron-atoms:1}
Escreva o \omterm{def:b3:many-electron-atoms:slater}{determinante de Slater} do estado fundamental do lítio, $1s^22s$,
com o elétron $2s$ de spin $\alpha$. Por que não há determinante para
$1s^3$?
\end{exercise}

\begin{exercise}[$\star$]\label{exo:b3:many-electron-atoms:2}
Conte os determinantes das configurações $p^2$, $p^3$ e $d^2$. Os termos
de $p^3$ são ${}^4$S, ${}^2$D, ${}^2$P e os de $d^2$ são ${}^3$F, ${}^3$P, ${}^1$G,
${}^1$D, ${}^1$S: verifique as duas contagens.
\end{exercise}

\begin{exercise}[$\star$]\label{exo:b3:many-electron-atoms:3}
Dê o termo e o nível fundamentais de N, O, F, \ce{Ti^2+} e \ce{Fe^3+} ($d^5$).
\end{exercise}

\begin{exercise}[$\star$]\label{exo:b3:many-electron-atoms:4}
Liste os níveis dos termos ${}^2$D e ${}^3$P com suas \omterm{def:b3:quantum-model-systems:degenerate}{degenerescências} $2J+1$
e verifique que as \omterm{def:b3:quantum-model-systems:degenerate}{degenerescências} somam $(2L+1)(2S+1)$.
\end{exercise}

\begin{exercise}[$\star\star$]\label{exo:b3:many-electron-atoms:5}
Os níveis ${}^3$P do carbono ficam em 0, 16.4 e \qty{43.4}{cm^{-1}}. Calcule a
razão dos intervalos e a constante spin--órbita $A$ a partir de cada intervalo.
O que diz a diferença entre os dois valores de $A$?
\end{exercise}

\begin{exercise}[$\star\star$]\label{exo:b3:many-electron-atoms:6}
Os níveis $3p$ do sódio ficam em 16956.17 e \qty{16973.37}{cm^{-1}}. Calcule
os comprimentos de onda no vácuo das raias D, sua separação em \unit{cm^{-1}} e
em meV, e $A$ para o termo $3p$.
\end{exercise}

\begin{exercise}[$\star\star$]\label{exo:b3:many-electron-atoms:7}
Quais destas transições do sódio são permitidas em emissão: $3p \to 3s$,
$3d \to 3s$, $3d \to 3p$, $4s \to 3s$, $4s \to 3p$? Justifique cada uma.
\end{exercise}

\begin{exercise}[$\star\star$]\label{exo:b3:many-electron-atoms:8}
% ledger: lev:He.2p3P2, lev:He.2p3P1, lev:He.2p3P0, lev:He.2p1P1
Os níveis $1s2p$ do hélio são \termsym{3}{P}{2}, \termsym{3}{P}{1},
\termsym{3}{P}{0} em 169086.76, 169086.84, \qty{169087.83}{cm^{-1}} e
\termsym{1}{P}{1} em \qty{171134.89}{cm^{-1}}. Calcule a energia média do
termo ${}^3$P e depois $K(1s,2p)$ em eV. Compare com $K(1s,2s) = \qty{0.398}{eV}$
e explique a diferença.
\end{exercise}

\begin{exercise}[$\star\star$]\label{exo:b3:many-electron-atoms:9}
Encontre os termos de $d^2$ pelo método deste capítulo, partindo do maior
$M_L$ (você pode dispensar a tabela completa, mas justifique cada termo).
\end{exercise}

\begin{exercise}[$\star\star\star$]\label{exo:b3:many-electron-atoms:10}
Mostre que a função espacial do tripleto $\frac{1}{\sqrt2}[1s(1)2s(2) -
2s(1)1s(2)]$ se anula quando $\mathbf r_1 = \mathbf r_2$ e que a do singleto
não se anula. Relacione isso com o sinal de $E_+ - E_-$.
\end{exercise}

\begin{exercise}[$\star\star\star$]\label{exo:b3:many-electron-atoms:11}
% ledger: lev:Ni2.3F3, lev:Ni2.3F2
Para \ce{Ni^2+} ($d^8$), os níveis ${}^3$F ficam em 0 ($J = 4$), 1360.7 ($J = 3$)
e \qty{2269.6}{cm^{-1}} ($J = 2$). Explique a ordem, calcule a razão dos
intervalos e compare com a previsão de Landé. Dê $A$ a partir do maior
intervalo.
\end{exercise}

\begin{exercise}[$\star\star\star$]\label{exo:b3:many-electron-atoms:12}
Mostre que o \omterm{def:b3:quantum-model-systems:hermitian}{valor esperado} de $\hat{\mathbf L}\cdot\hat{\mathbf S}$, somado
sobre todos os $(2L+1)(2S+1)$ estados de um termo, é nulo, de modo que o
\omterm{def:b3:many-electron-atoms:spin-orbit}{acoplamento spin--órbita} deixa inalterada a média ponderada dos níveis.
Verifique no ${}^3$P do carbono.
\end{exercise}

\section{Problema: Dois hélios}

\begin{problem}[{Problema de fim de semana --- as famílias singleto e tripleto
do hélio: determinantes, os termos dos íons de metais de transição, as regras
de seleção que separam as duas famílias e a integral de troca medida}]
\label{pb:b3:many-electron-atoms:1}
% ledger: lev:He.2s3S1, lev:He.2s1S0, lev:He.2p3P2, lev:He.2p3P1, lev:He.2p3P0, lev:He.2p1P1, lev:Ti2.3F3, lev:Ti2.3F4, lev:Ti2.3P0, lev:Ti2.3P1, lev:Ti2.3P2
Dados (NIST, em \unit{cm^{-1}} acima do estado fundamental de cada espécie). Hélio:
$1s2s$ \termsym{3}{S}{1} 159855.97, \termsym{1}{S}{0} 166277.44; $1s2p$
\termsym{3}{P}{2,1,0} 169086.76, 169086.84, 169087.83, \termsym{1}{P}{1}
171134.89. \ce{Ti^2+}: \termsym{3}{F}{2,3,4} 0, 184.9, 420.4; \termsym{3}{P}{0,1,2}
10538.4, 10603.6, 10721.2. $\qty{1}{eV} = \qty{8065.54}{cm^{-1}}$.

\medskip\noindent\textbf{Parte I --- \omterm{def:b3:many-electron-atoms:spin-orbital}{Spin-orbitais} e determinantes.}
\begin{enumerate}
  \item Liste os quatro \omterm{def:b3:many-electron-atoms:spin-orbital}{spin-orbitais} disponíveis para a configuração $1s2s$.
  \item Escreva os quatro determinantes com um elétron em $1s$ e um em
        $2s$.
  \item Mostre que $|1s\alpha\,2s\alpha|$ é o produto de uma função espacial
        antissimétrica pela função de spin $\alpha(1)\alpha(2)$.
  \item Combine $|1s\alpha\,2s\beta|$ e $|1s\beta\,2s\alpha|$ para obter a
        componente $M_S = 0$ da mesma função espacial, e um quarto estado.
  \item Qual função espacial, simétrica ou antissimétrica, acompanha
        $S = 0$? E $S = 1$?
  \item Por que eles se chamam singleto e tripleto?
\end{enumerate}

\medskip\noindent\textbf{Parte II --- Os termos de \ce{Ti^2+}.}
\begin{enumerate}[resume]
  \item Dê a configuração de \ce{Ti^2+} e o número de seus
        determinantes.
  \item Encontre seu termo e seu nível fundamentais com as regras de Hund;
        confira com os dados.
  \item Seus outros termos são ${}^3$P, ${}^1$G, ${}^1$D e ${}^1$S: verifique a
        contagem de estados.
  \item A \omterm{def:b3:many-electron-atoms:spin-orbit}{estrutura fina} de ${}^3$F é normal ou invertida? Teste a regra dos
        intervalos de Landé.
  \item Calcule as energias médias ponderadas dos termos ${}^3$F e ${}^3$P.
  \item Sua separação é $15B$, onde $B$ é um parâmetro de Racah do
        íon (\cref{ch:b3:complex-spectra-magnetism}). Calcule $B$.
\end{enumerate}

\medskip\noindent\textbf{Parte III --- Por que duas famílias.}
\begin{enumerate}[resume]
  \item Qual regra de seleção proíbe uma transição entre um singleto e um
        tripleto?
  \item O tripleto mais baixo, $1s2s$ \termsym{3}{S}{1}, não pode emitir para o
        estado fundamental $1s^2$ \termsym{1}{S}{0}. Dê duas razões. Como se
        chama um estado assim?
  \item Calcule o comprimento de onda da raia $1s2p\ {}^3$P $\to 1s2s\ {}^3$S (use
        a média dos níveis ${}^3$P).
  \item Calcule o comprimento de onda de $1s2p\ {}^1$P $\to 1s2s\ {}^1$S.
  \item Calcule o comprimento de onda de $1s2p\ {}^1$P $\to 1s^2\ {}^1$S. Em que
        região do espectro ele está?
  \item Explique por que os espectroscopistas do século XIX, vendo essas raias,
        acreditavam em dois gases diferentes.
\end{enumerate}

\medskip\noindent\textbf{Parte IV --- A \omterm{def:b3:many-electron-atoms:exchange}{integral de troca}, medida.}
\begin{enumerate}[resume]
  \item Escreva as energias de primeira ordem do singleto e do tripleto $1s2s$
        em função de $E_{1s}$, $E_{2s}$, $J$ e $K$.
  \item Qual é o mais baixo, e por quê, em termos das posições dos elétrons?
  \item Calcule a separação singleto--tripleto $1s2s$ em \unit{cm^{-1}} e em eV.
  \item Calcule $K(1s,2p)$ em eV a partir dos níveis $1s2p$.
  \item Por que $K(1s,2s)$ é maior que $K(1s,2p)$?
  \item A primeira regra de Hund é confirmada pelos estados excitados do hélio?
  \item Enuncie o resultado: a \omterm{def:b3:many-electron-atoms:exchange}{integral de troca} $K(1s,2s)$ do hélio, em
        eV.
\end{enumerate}
\end{problem}
